% Options for packages loaded elsewhere
\PassOptionsToPackage{unicode}{hyperref}
\PassOptionsToPackage{hyphens}{url}
\documentclass[
  ignorenonframetext,
  aspectratio=169,
]{beamer}
\newif\ifbibliography
\usepackage{pgfpages}
\setbeamertemplate{caption}[numbered]
\setbeamertemplate{caption label separator}{: }
\setbeamercolor{caption name}{fg=normal text.fg}
\beamertemplatenavigationsymbolsempty
% remove section numbering
\setbeamertemplate{part page}{
  \centering
  \begin{beamercolorbox}[sep=16pt,center]{part title}
    \usebeamerfont{part title}\insertpart\par
  \end{beamercolorbox}
}
\setbeamertemplate{section page}{
  \centering
  \begin{beamercolorbox}[sep=12pt,center]{section title}
    \usebeamerfont{section title}\insertsection\par
  \end{beamercolorbox}
}
\setbeamertemplate{subsection page}{
  \centering
  \begin{beamercolorbox}[sep=8pt,center]{subsection title}
    \usebeamerfont{subsection title}\insertsubsection\par
  \end{beamercolorbox}
}
% Prevent slide breaks in the middle of a paragraph
\widowpenalties 1 10000
\raggedbottom
\AtBeginPart{
  \frame{\partpage}
}
\AtBeginSection{
  \ifbibliography
  \else
    \frame{\sectionpage}
  \fi
}
\AtBeginSubsection{
  \frame{\subsectionpage}
}
\usepackage{iftex}
\ifPDFTeX
  \usepackage[T1]{fontenc}
  \usepackage[utf8]{inputenc}
  \usepackage{textcomp} % provide euro and other symbols
\else % if luatex or xetex
  \usepackage{unicode-math} % this also loads fontspec
  \defaultfontfeatures{Scale=MatchLowercase}
  \defaultfontfeatures[\rmfamily]{Ligatures=TeX,Scale=1}
\fi
\usepackage{lmodern}
\ifPDFTeX\else
  % xetex/luatex font selection
\fi
% Use upquote if available, for straight quotes in verbatim environments
\IfFileExists{upquote.sty}{\usepackage{upquote}}{}
\IfFileExists{microtype.sty}{% use microtype if available
  \usepackage[]{microtype}
  \UseMicrotypeSet[protrusion]{basicmath} % disable protrusion for tt fonts
}{}
\makeatletter
\@ifundefined{KOMAClassName}{% if non-KOMA class
  \IfFileExists{parskip.sty}{%
    \usepackage{parskip}
  }{% else
    \setlength{\parindent}{0pt}
    \setlength{\parskip}{6pt plus 2pt minus 1pt}}
}{% if KOMA class
  \KOMAoptions{parskip=half}}
\makeatother
\usepackage{longtable,booktabs,array}
\usepackage{caption}
\captionsetup[table]{skip=6pt}
\newcounter{none} % for unnumbered tables
\usepackage{calc} % for calculating minipage widths
\usepackage{caption}
% Make caption package work with longtable
\makeatletter
\def\fnum@table{\tablename~\thetable}
\makeatother
\setlength{\emergencystretch}{3em} % prevent overfull lines
\providecommand{\tightlist}{%
  \setlength{\itemsep}{0pt}\setlength{\parskip}{0pt}}
% Auto-generated by infrastructure/rendering/_pdf_latex_helpers.py::extract_math_font_preamble.
% Minimal header passed to Pandoc via -H so Beamer slides pick up the same math font
% as the combined-PDF path. Adjust the manuscript's preamble.md to change the font globally.
\usepackage{unicode-math}
\setmathfont{latinmodern-math.otf}

% Manuscript macro/package fallbacks (newcommand -> providecommand).
% Core mathematics
\usepackage{amsmath}
\usepackage{amssymb}
% Document layout
% Tables
\usepackage{booktabs}
% Caption styling (booktabs already loaded above). Small caption font with a
% bold label ("Figure 1", "Table 2") gives the math/figure-heavy layout a
% consistent, journal-like caption rhythm without fighting pandoc-crossref —
% pandoc emits standard \caption{}, and caption/captionsetup only restyle it.
% ── Theorem environments ─────────────────────────────────────────────────
% amsthm is loaded above. The FedGVI<->Friston bridge is stated as a small
% number of theorems/definitions/lemmas/propositions/corollaries; sharing one
% counter (definition/lemma/proposition/corollary all step the `theorem`
% counter) gives a single monotone numbering 1, 2, 3, ... across all kinds, so
% authors can reference them as Theorem 1, Definition 2, Corollary 3 without a
% per-kind counter drifting out of order. Plain style for theorem-like results,
% definition style (upright body) for definitions.
% (algorithm2e intentionally not loaded: this manuscript states results as
% numbered amsthm Theorems/Definitions, not algorithm2e pseudocode.)
% Code listings
\usepackage{listings}
% Typography and formatting
% (siunitx intentionally not loaded: no \SI/\num units appear in this manuscript.)
% Cross-references and citations
% (cleveref and titlesec intentionally not loaded: unavailable in this TeX tree
% and not required. Cross-references resolve through pandoc-crossref's own
% \ref-style names — no \cref appears — and section heading styling falls back
% to the pandoc/LaTeX defaults.)
% ── TeX Live Basic-compatible mono font for code listings ────────────
% Latin Modern Mono is bundled with TeX Live Basic and keeps the manuscript
% renderable on minimal TeX installations. Mathematical Unicode coverage is
% provided separately by Latin Modern Math below, so the code font does not
% need a private JuliaMono installation.
%
% Use the TeX font filename rather than the system-family name: the minimal
% TeX Live fontconfig database may not register the OpenType family even though
% kpsewhich can resolve the bundled files.
% Math font for unicode-math: Latin Modern Math (TeX Live) has full BMP
% coverage including U+2223 (\mid), U+226A/226B (\ll/\gg), and the Greek/
% blackboard letters used in equations. Without an explicit \setmathfont,
% unicode-math falls back to lmroman text font which lacks several glyphs
% and emits "Missing character" warnings on every \mid in math mode.

% Slide layout defaults for warning-clean scientific decks.
\usepackage{etoolbox}
\IfFileExists{xurl.sty}{\usepackage{xurl}}{}
\IfFileExists{seqsplit.sty}{\usepackage{seqsplit}}{\newcommand{\seqsplit}[1]{#1}}
\protected\def\breakseq#1{\seqsplit{#1}}
\protected\def\breaktt#1{\begingroup\ttfamily\seqsplit{#1}\endgroup}
\setlength{\emergencystretch}{6em}
\tolerance=5000
\hbadness=10000
\hfuzz=1pt
\setlength{\tabcolsep}{2pt}
\AtBeginEnvironment{longtable}{\tiny\renewcommand{\arraystretch}{0.86}\setlength{\tabcolsep}{1pt}}
\AtBeginEnvironment{tabular}{\tiny\renewcommand{\arraystretch}{0.86}\setlength{\tabcolsep}{1pt}}
\AtBeginEnvironment{equation}{\tiny}
\AtBeginEnvironment{equation*}{\tiny}
\AtBeginEnvironment{align}{\tiny}
\AtBeginEnvironment{align*}{\tiny}
\AtBeginEnvironment{itemize}{\footnotesize}
\AtBeginEnvironment{enumerate}{\footnotesize}
\AtBeginEnvironment{description}{\footnotesize}
\setbeamerfont{caption}{size=\tiny}
\setbeamerfont{caption name}{size=\tiny}
\setbeamerfont{normal text}{size=\small}
\setbeamerfont{frametitle}{size=\small}
\setbeamerfont{section title}{size=\footnotesize}
\setbeamerfont{subsection title}{size=\footnotesize}
\setbeamertemplate{section page}{%
  \centering
  \begin{beamercolorbox}[sep=12pt,center,wd=\paperwidth]{section title}
    \parbox{0.86\paperwidth}{\centering\usebeamerfont{section title}\insertsection\par}
  \end{beamercolorbox}
}
\setbeamertemplate{subsection page}{%
  \centering
  \begin{beamercolorbox}[sep=8pt,center,wd=\paperwidth]{subsection title}
    \parbox{0.86\paperwidth}{\centering\usebeamerfont{subsection title}\insertsubsection\par}
  \end{beamercolorbox}
}
\setlength{\abovecaptionskip}{2pt}
\setlength{\belowcaptionskip}{0pt}

% Natbib and cross-reference fallbacks — slides don't load natbib
% or cleveref, but combined-PDF manuscript prose may emit these
% commands. The fallback renders citations as a bracketed key list
% and cross-references as detokenized labels so slides don't fail on
% undefined control sequences or raw underscores. \providecommand is
% a no-op if packages load later.
\providecommand{\citep}[1]{[#1]}
\providecommand{\citet}[1]{#1}
\providecommand{\citealp}[1]{#1}
\providecommand{\citeauthor}[1]{#1}
\providecommand{\citeyear}[1]{#1}
\providecommand{\cref}[1]{\texttt{\detokenize{#1}}}
\providecommand{\Cref}[1]{\texttt{\detokenize{#1}}}
\providecommand{\crefrange}[2]{\texttt{\detokenize{#1}}--\texttt{\detokenize{#2}}}
\providecommand{\Crefrange}[2]{\texttt{\detokenize{#1}}--\texttt{\detokenize{#2}}}
\providecommand{\paragraph}[1]{\textbf{#1}\ }

% Opt-in accessible presentation profile.
\setbeamerfont{normal text}{size*={20pt}{24pt}}
\setbeamerfont{frametitle}{size*={28pt}{32pt}}
\setbeamerfont{section title}{size*={28pt}{32pt}}
\setbeamerfont{subsection title}{size*={28pt}{32pt}}
\setbeamerfont{caption}{size*={16pt}{19pt}}
\setbeamerfont{caption name}{size*={16pt}{19pt}}
\setbeamerfont{itemize/enumerate body}{size*={20pt}{24pt}}
\setbeamerfont{itemize/enumerate subbody}{size*={20pt}{24pt}}
\setbeamerfont{itemize/enumerate subsubbody}{size*={20pt}{24pt}}
\setbeamerfont{itemize item}{size*={20pt}{24pt}}
\setbeamerfont{itemize subitem}{size*={20pt}{24pt}}
\setbeamerfont{itemize subsubitem}{size*={20pt}{24pt}}
\setbeamerfont{description body}{size*={20pt}{24pt}}
\setbeamerfont{description item}{size*={20pt}{24pt}}
\setbeamerfont{quote}{size*={20pt}{24pt}}
\setbeamerfont{quotation}{size*={20pt}{24pt}}
\AtBeginEnvironment{longtable}{\fontsize{20pt}{24pt}\selectfont}
\AtBeginEnvironment{tabular}{\fontsize{20pt}{24pt}\selectfont}
\AtBeginEnvironment{equation}{\fontsize{20pt}{24pt}\selectfont}
\AtBeginEnvironment{equation*}{\fontsize{20pt}{24pt}\selectfont}
\AtBeginEnvironment{align}{\fontsize{20pt}{24pt}\selectfont}
\AtBeginEnvironment{align*}{\fontsize{20pt}{24pt}\selectfont}
\AtBeginEnvironment{itemize}{\fontsize{20pt}{24pt}\selectfont}
\AtBeginEnvironment{enumerate}{\fontsize{20pt}{24pt}\selectfont}
\AtBeginEnvironment{description}{\fontsize{20pt}{24pt}\selectfont}
\AtBeginDocument{\usebeamerfont{normal text}}
\newlength{\AFSlideTableWidth}
% The fixed footer reservation is calibrated with the semantic seven-line regular-body budget.
\setbeamertemplate{footline}{%
  \leavevmode\vbox{%
    \hbox{%
      \begin{beamercolorbox}[wd=\paperwidth,ht=16pt,dp=3pt,center]{author in head/foot}%
        \fontsize{16pt}{19pt}\selectfont Untagged PDF derivative \textbar\ \href{../web/index.html}{HTML reader}%
      \end{beamercolorbox}%
    }%
    \vskip6pt%
  }%
}
% Accessible footer reserved height: 25pt.

\newtheorem{proposition}{Proposition}
\newtheorem{hypothesis}{Hypothesis}
\newtheorem{remark}[theorem]{Remark}
\newtheorem{axiom}{Axiom}
\newtheorem{property}{Property}
\usepackage{bookmark}
\IfFileExists{xurl.sty}{\usepackage{xurl}}{} % add URL line breaks if available
\urlstyle{same}
% fallback for those not using the hyperref driver hyperxmp:
\makeatletter
\@ifundefined{xmpquote}{\newcommand{\xmpquote}[1]{#1}}{}
\makeatother
\hypersetup{
  hidelinks,
  pdfcreator={LaTeX via pandoc}}

\author{\texorpdfstring{}{}}
\date{}

\begin{document}

\section{Methods: the federated active-inference
stack}\label{sec:methods}

\begin{frame}{Methods: the federated active-inference stack}
\protect\phantomsection\label{sec:methodology}{}

This section develops the federated generalized variational inference
(FedGVI) core in the discrete-categorical setting and defines its
primitives. Their recovery limits are stated as numbered theorems in
sec.~6.
\end{frame}

\begin{frame}{Methods: the federated\ldots{} (part 2)}
Every belief here is a categorical pmf --- a non-negative vector summing
to one --- so the generalized-variational-inference machinery reduces to
closed forms that are exactly testable.
\end{frame}

\begin{frame}[fragile]{Methods: the federated\ldots{} (part 3)}
All mathematics lives in \texttt{src/fedference/}; the prose names the
module and the identity that pins each claim.
\end{frame}

\begin{frame}{Methods: the federated\ldots{} (part 4)}
The active-inference community has built a rich apparatus for federated
belief sharing: discrete-state-space agents that broadcast posteriors
and fuse them into a consensus (Da Costa et al. 2020; Friston et al.
2024),
\end{frame}

\begin{frame}{Methods: the federated\ldots{} (part 5)}
message-passing toolboxes that make exact-Bayes inference scalable
(Heins et al. 2022; Bagaev et al. 2023), and collective and multi-agent
formulations in which ensembles coordinate by sharing observations and
beliefs (Heins et al. 2024; Albarracin et al. 2022; Kaufmann et al.
2021).
\end{frame}

\begin{frame}{Methods: the federated\ldots{} (part 6)}
We accept that apparatus and extend it. Among the active-inference
sources reviewed here, the consensus rules are exact-Bayes and trusting,
without an account of an ensemble member that is misspecified or
adversarial.
\end{frame}

\begin{frame}{Methods: the federated\ldots{} (part 7)}
Outside active inference, the federated-learning and robust-Bayes
literatures address important parts of that question --- decentralized
aggregation (McMahan et al. 2017), partitioned and federated variational
inference (Ashman et al. 2022; Bui et al. 2018), and generalized,
\end{frame}

\begin{frame}{Methods: the federated\ldots{} (part 8)}
robustness-bearing Bayesian updating (Bissiri et al. 2016; Knoblauch et
al. 2022; Basu et al. 1998; Zhang and Sabuncu 2018). Among the sources
reviewed here, that apparatus has not been carried into the
generative-model-bearing, action-selecting POMDP setting. The
methodology below is the bridge.
\end{frame}

\begin{frame}{Methods: the federated\ldots{} (part 9)}
It federates the FedGVI objective (Mildner et al. 2025) per agent inside
an active-inference ensemble, proves the standard-Bayes client limits,
and tests the project-local zero-robustness log-linear-pool identity.
Under the qualified categorical bridge of
sec.~5.8,
\end{frame}

\begin{frame}{Methods: the federated\ldots{} (part 10)}
that pool specializes Eq. 7's message-combination term rather than the
complete source protocol. fig.~2 illustrates the
three-axis architecture and the recovery hierarchy.
\end{frame}

\begin{frame}{Federation protocol: local update, server fusion,
broadcast}
\protect\phantomsection\label{sec:method-protocol}
A colony of \(N\) agents shares a single latent factor
\(s \in \{1,\dots,n_s\}\) (in the sentinel scenario, the location of a
creature on a grid of \(n_s\) cells). Each round proceeds in three
steps:
\end{frame}

\begin{frame}{1. Local inference}
\protect\phantomsection\label{local-inference}
Agent \(n\) observes \(o_n\) and forms a local posterior \(q_n(s)\) over
the shared factor by a generalized-Bayes update against its own cavity,
the colony belief with agent \(n\)'s previous contribution removed.
\end{frame}

\begin{frame}{1. Local inference (part 2)}
This is where robustness enters per agent: the update minimizes a
loss-plus-divergence objective, and the FedGVI choice of a bounded loss
carries the source theorem's bounded-influence result under its matching
assumptions.
\end{frame}

\begin{frame}{2. Broadcast}
\protect\phantomsection\label{broadcast}
Agent \(n\) broadcasts \(q_n(s)\), optionally with a scalar base weight
\(w_n \ge 0\).
\end{frame}

\begin{frame}{3. Aggregation}
\protect\phantomsection\label{aggregation}
The server, or equivalently each agent acting as its own server, fuses
the broadcast beliefs into a consensus. Following sensory attenuation
--- ``agents do not hear themselves'' --- an agent's heard consensus
excludes its own message.
\end{frame}

\begin{frame}{3. Aggregation (part 2)}
The protocol has two distinct places where robustness can live, and we
keep them separate throughout. The \textbf{per-agent generalized-Bayes
update} in step 1 is FedGVI-faithful at the stated primitive level: its
formal bounded-influence claim is conditional on the source theorem's
assumptions.
\end{frame}

\begin{frame}{3. Aggregation (part 3)}
The \textbf{server-side aggregation rule} in step 3 admits an optional
divergence-reweighting heuristic that down-weights agents far from the
emerging consensus; this heuristic is a complementary device whose
positive formal property is recovery of the naive consensus in its
trusting limit,
\end{frame}

\begin{frame}{3. Aggregation (part 4)}
while a scoped proposition rejects one declared separable objective
class.

sec.~3.3 holds this boundary; no figure, table, or
sentence in this work grants the server-side heuristic the per-agent
FedGVI guarantee.
\end{frame}

\begin{frame}{Notation for beliefs, losses, and divergences}
\protect\phantomsection\label{sec:method-notation}
The authoritative symbol and API contract is
sec.~14. In the main text, \(q_n(s)\) denotes agent
\(n\)'s local posterior, \(q(s)\) the global consensus, and
\(q_{-n}(s)\) the cavity after removing the site factor \(t_n(s)\).
\end{frame}

\begin{frame}{Notation for beliefs, losses\ldots{} (part 2)}
The prior is \(\pi_0(s)\), while \(\boldsymbol{\pi}\) is a policy. The
POMDP tensors are \(A[o,s]\), \(B[s',s,u]\), \(C[o]\), and \(D_0[s]\).
\end{frame}

\begin{frame}{Notation for beliefs, losses\ldots{} (part 3)}
The aggregation weights are \(w_n\) (raw/base), \(a_n\) (raw variational
effective), and \(\widetilde a_n\) (normalized influence). The server
robustness coefficient is \(c\), the variational entropy weight is
\(\lambda\), the Rényi order is \(\alpha\), the density-power parameter
is \(\beta\),
\end{frame}

\begin{frame}{Notation for beliefs, losses\ldots{} (part 4)}
and the robust cross-entropy parameter is \(q_{\text{loss}}\). The
notation supplement also defines the seed/trial nesting and all
statistical quantities used below.
\end{frame}

\begin{frame}{Notation for beliefs, losses\ldots{} (part 5)}
The study is run over a fixed ensemble of 7 agents sharing a factor of 9
locations, with all randomness seeded at 0; the full per-study
configuration is tabulated in tbl.~2. As an
independent generative-model-free baseline,
\end{frame}

\begin{frame}{Notation for beliefs, losses\ldots{} (part 6)}
we also implement FedGVI in a deterministic MLP complement trained with
the density-power \(\beta\)-loss --- generalized variational inference
with a point-mass variational family (sec.~7.6).
\end{frame}

\begin{frame}{Notation for beliefs, losses\ldots{} (part 7)}
The remaining methodology subsections develop each primitive in turn:
the generalized-Bayes update and its recovery to standard Bayes
(sec.~5.3),
\end{frame}

\begin{frame}{Notation for beliefs, losses\ldots{} (part 8)}
the divergence family and its KL limit
(sec.~5.6) and the robust loss family and its
NLL limit (sec.~5.7), the aggregation identity
(sec.~5.8), the lift to a belief-sharing round
(sec.~5.9),
\end{frame}

\begin{frame}{Notation for beliefs, losses\ldots{} (part 9)}
and the paired statistics that earn every ``robust beats naive'' verdict
(sec.~5.27).
\end{frame}

\end{document}
